
\chapter{Процессинговый центр}\label{ch:processing-center}
Процессинговый центр связывает банк и~карточную сеть: принимает сообщение
ISO~8583, выбирает маршрут, выполняет криптографические проверки, ведёт журнал
операций и~позже превращает авторизацию в~клиринговую и~расчётную запись.

Процессинг существует в~трёх конфигурациях. \textbf{Acquirer processor}
обслуживает эквайера: приём сообщений от~терминалов и~шлюзов, маршрутизация
авторизаций в~сеть, формирование исходящего клирингового файла, разбор входящих
расчётных сообщений. \textbf{Issuer processor} обслуживает эмитента: CMS,
верификация ARQC и~генерация ARPC, обработка входящих клиринговых записей
по~своему BIN, issuer scripting, токенизация на~стороне эмитента.
\textbf{Coupled processor} объединяет обе стороны и~обрабатывает on-us-операции
напрямую, минуя сеть.

По~Федеральному закону от~27.06.2011~№~161-ФЗ \enquote{О~национальной
платёжной системе} процессор отдельного банка остаётся его внутренней
инфраструктурой~\cite{fz-161}. Операционным и~платёжным клиринговым центром
(ОПКЦ) платёжной системы выступает выделенный оператор: НСПК~--- для~ПС~Мир
и~по~поручению Банка России для~СБП (оператор СБП~--- сам Банк России).
Самостоятельный процессор-вендор, обслуживающий несколько участников, при
выполнении критериев включается в~реестр операторов услуг платёжной
инфраструктуры (ОУПИ) Банка России в~отдельном порядке
(см.~раздел~\ref{sec:pc-licensing}).

Внутри процессинг состоит из~коммутатора авторизации, пула HSM, клирингового
и~расчётного модулей; каждый компонент задаёт свои требования к~задержке,
доступности и~контролю над ключевым материалом.
Распределение этих компонентов по~трём конфигурациям сведено
на~рис.~\ref{fig:processor-configurations}.

\begin{figure}[tbp]
  \centering
  \includefiguresvg[width=\linewidth]{assets/figures/ch36-processing-center/processor-configurations}
  \caption{Распределение функций процессингового центра
  по~конфигурациям: acquirer, issuer, coupled.}\label{fig:processor-configurations}
\end{figure}

\section{Компоненты процессингового центра}\label{sec:pc-components}

\subsection*{Коммутатор авторизации}

Коммутатор авторизации принимает каждое входящее сообщение ISO~8583
от~эквайеров или от~шлюзов, действующих от~их имени, выбирает маршрут
к~эмитенту или в~карточную сеть и~возвращает ответ. Он работает в~реальном
времени, по~одному сообщению: допустимая задержка измеряется десятками
миллисекунд.~\cite{visa-core-rules,mc-tpr}

Таблица BIN задаёт маршрут внутри коммутатора: для каждого диапазона BIN в~ней
указан канал отправки запроса~--- VisaNet, Banknet, ОПКЦ НСПК или прямое
подключение к~эмитенту. Карточные сети регулярно публикуют изменения
к~таблице BIN, и~процессинг обязан применять их в~установленные сроки. Иначе
часть транзакций будет отправлена по~неверному маршруту и~получит отказ
не~по~вине клиента, ТСП или эквайера.~\cite{visa-core-rules,mc-rules}

Коммутатор извлекает BIN (первые 6--8 цифр PAN из~DE~2) и~ищет совпадение
в~таблице BIN, определяя карточную сеть~--- Visa, Mastercard или Мир. Далее
проверяет, принадлежит ли BIN тому же банку, в~котором работает процессинг.
Такой случай называется on-us: сообщение маршрутизируется напрямую
к~авторизационному хосту эмитента, минуя сеть, а~interchange не начисляется.
Для off-us коммутатор выбирает исходящий канал (VisaNet, Banknet или
ОПКЦ), форматирует сообщение по~Interface Specification (IFS) конкретной сети
и~отправляет запрос. Различия между сетями проявляются в~DE~25, DE~48
и~в~структуре подэлементов (subelement layout).

Во время маршрутизации коммутатор проверяет обязательные поля, формат MTI
и~корректность bitmap, дополняет запрос данными, требуемыми сетью, и~ведёт
журнал транзакций. Запись с~точностью до~миллисекунд используют для отладки
и~последующей сверки с~клиринговыми файлами.

Когда внешний маршрут недоступен, включается резервная авторизация
\textit{stand-in} (механизм описан в~разделе~\ref{sec:tx-authorization}).
Помимо сетевого \textit{stand-in}, issuer processor может реализовать локальный
\textit{stand-in} по~поручению эмитента, опираясь на~собственные правила: лимит
суммы, диапазон BIN, историю операций по~карте. \textit{Stand-in} ограничивают
малыми суммами и~заранее определёнными категориями транзакций, потому что
одобрение проходит без онлайн-проверки баланса
и~лимитов.~\cite{visa-core-rules,mc-rules}

Распределение кредитного риска в~\textit{stand-in} следует из~того, кто его
выполнял. Если \textit{stand-in} выполняла карточная сеть (VisaNet
\textit{stand-in}, Mastercard \textit{stand-in}), сеть берёт на~себя кредитный
риск по~одобренным операциям в~пределах параметров \textit{stand-in}: лимита
суммы, списка BIN, ограничений MCC. Если \textit{stand-in} выполнял локальный
issuer processor по~поручению эмитента, кредитный риск остаётся на~эмитенте~---
именно он отвечает перед платёжной системой и~держателем карты, а~договор между
эмитентом и~процессором может предусматривать операционную ответственность
процессора за~превышение согласованных параметров \textit{stand-in}. Параметры
\textit{stand-in} (максимальная сумма, разрешённые MCC, диапазоны BIN) задают
предел этой операционной ответственности.

\subsection*{Пул HSM}

Устройство HSM, иерархия ключей, key ceremony и~базовые операции (проверка ARQC
  (\textit{Authorization Request Cryptogram}), генерация ARPC
(\textit{Authorization Response Cryptogram}), PIN translation) разобраны
в~главе~\ref{ch:crypto}.

Криптографические операции нельзя распределять по~узлам так же свободно, как
операции stateless-сервисов: ключевой материал загружается через
регламентированные процедуры, и~каждое новое устройство требует отдельной key
ceremony. Процессинговые центры строят пул HSM с~балансировкой нагрузки
и~заранее планируют его ёмкость. Производительность (число
переводов PIN, подписей RSA и~проверок ARQC в~секунду) зависит от~модели,
версии HSM и~лицензированных функций. Её сверяют по~документации
производителя~\cite{pci-pin-security,pci-pts-hsm,ansi-x9-24}. На российском
рынке кроме зарубежных моделей Thales payShield, Utimaco и~Atalla есть
отдельный класс сертифицированных ФСБ России средств криптографической защиты
информации (СКЗИ): от~КриптоПро HSM до~решений Аладдин Р.Д.\ и~ИнфоТеКС. Без
них эксплуатация криптографии по~ГОСТ невозможна. При росте нагрузки узким
местом становится HSM; коммутатор и~сеть масштабируются проще.

\importantnote{%
  Горизонтальное масштабирование пула HSM ограничено физически:
  ввод каждого нового устройства требует key ceremony
  (см.~раздел~\ref{sec:crypto-key-ceremony}). Планирование ёмкости
  учитывает текущую нагрузку и~часы, необходимые на~ввод устройства
  в~эксплуатацию.
}

Синхронизация KCV (\textit{Key Check Value}) между устройствами подтверждает,
что все HSM пула работают с~одним и~тем же ключом, не~раскрывая самого ключа.
При отказе одного HSM ключевой материал не~теряется: он загружен во~все
устройства пула в~одной key ceremony. Нарастить ёмкость за~минуты при этом
невозможно~--- церемония занимает часы.

\subsection*{Клиринговый модуль}

Клиринговый модуль, в~отличие от~авторизации, обрабатывает пакеты.

На~стороне acquirer processor модуль запускается по~расписанию, собирает
авторизованные транзакции за~период и~добавляет данные, которых на~этапе
авторизации ещё не~было: финальную сумму после добавления чаевых, результат
адресной проверки, параметры конверсии DCC. Затем формирует исходящий
клиринговый файл в~формате сети: TC33 (\textit{transaction code 33},
capture-формат Visa в~BASE~II) для~Visa, IPM (\textit{Integrated Product
Messages}, клиринговый формат Mastercard на~MTI~1240) для~Mastercard или
профиль ISO~8583 с~расширениями НСПК для~внутрироссийских операций (доставка
  через систему электронного документооборота НСПК;
см.~раздел~\ref{sec:recon-mir}).~\cite{visa-core-rules,mc-tpr,nspk-mir-rules}

На~стороне issuer processor модуль принимает входящие клиринговые файлы от~сети
и~сопоставляет каждую запись с~ранее одобренной авторизацией по~полям STAN,
RRN, сумме и~дате. Coupled processor видит обе стороны и~сопоставляет свой
исходящий клиринг со~своим входящим внутри одной системы.

В~любой конфигурации сопоставление вскрывает расхождения: \textit{orphan
authorization} (авторизация без клиринга), \textit{late presentment} (поздний
клиринг), \textit{force post} (клиринг без авторизации), частичное
подтверждение и~\textit{tip adjustment}~--- корректировка клиринговой суммы
на~величину чаевых, добавленных после авторизации. Любое такое расхождение
создаёт риск прямой денежной потери; клиринговый модуль формирует отчёт
об~исключениях, который команда сопровождения разбирает по~правилам сверки
(см.~раздел~\ref{sec:recon-gaps}).

\subsection*{Расчётный модуль}

Нетто-позиции всех участников взаимозачёта по~операциям конкретной сети считает
сама сеть (BASE~II у~Visa, GCMS у~Mastercard, ОПКЦ НСПК для~внутрироссийских
операций), не~процессор отдельного банка. Сеть рассылает каждому банку
расчётные сообщения с~его собственной нетто-позицией и~детализацией:
interchange, комиссии сети, корректировки (чарджбэки, повторные представления,
сборы).

Расчётный модуль процессора принимает эти сообщения и~формирует на~их основе
проводки во~внутренней банковской системе. Для внутрироссийских транзакций
расчёт проходит через НСПК и~платёжную систему Банка России
(см.~раздел~\ref{sec:banks-infrastructure}), для~международных сценариев~---
через расчётные механизмы, предусмотренные правилами соответствующей
схемы.~\cite{cbr-ps,visa-core-rules,mc-switching-services} На расчётном этапе
деньги уже движутся между банками, и~любая ошибка в~проводках превращается
в~прямую финансовую потерю.

\section{Производительность и~отказоустойчивость}\label{sec:pc-performance}

Требования к~задержке и~доступности процессингового центра задаются внешним
контрактом сети. Visa публикует maximum time limits for Authorization Request
Response и~связывает просрочку ответа со~включением \textit{stand-in};
Mastercard в~Transaction Processing Rules задаёт authorization performance
standards для участников сети.~\cite{visa-core-rules,mc-tpr}

\subsection*{Бюджет задержки}

Карточные сети нормируют внешний предел времени ответа, но~не дают
универсальной внутренней таблицы \enquote{сколько миллисекунд должен занимать
каждый сервис}. Внутренний бюджет обязан оставлять запас на~сетевой путь,
обработку у~эмитента и~возврат ответа. Пример внутреннего бюджета: парсинг
и~валидация сообщения 1--3~мс, обращение к~HSM (ARQC или
перевод PIN) 3--7~мс, поиск по~BIN и~маршрутизация 1~мс, логирование 1--2~мс,
сериализация ответа 1~мс. Внутренняя обработка укладывается в~7--15~мс;
остальной бюджет остаётся на~сетевой путь до~эмитента и~обратно.

\subsection*{Резервирование и~геораспределение}

Несколько минут недоступности процессингового центра превращаются в~тысячи
отклонённых транзакций и~штрафы от~карточных сетей. Правила сетей
не~предписывают единственную допустимую схему (active-active, active-standby
или гибридную), но~в~любом варианте оператор должен выдерживать отказ площадки
без потери управления трафиком.

Active-active предпочтительнее active-standby из-за времени запуска пассивной
площадки. При переключении она переходит от~холодного состояния к~полной
нагрузке: HSM загружает ключи, таблицы BIN синхронизируются, очереди
восстанавливаются. Каждая секунда запуска означает отказанные транзакции.
В~active-active оба узла работают под~реальной нагрузкой, и~при отказе одного
трафик перераспределяется без отдельного запуска резервной площадки. Цена
решения~--- риск split-brain при~разрыве репликации. Авторизация не~хранит
состояния между запросами, и~последствия кратковременного split-brain для неё
ограничены. Клирингу и~расчёту нужна более жёсткая гарантия согласованности:
авторизационный след досверяется по~внешним журналам сети, а~потеря клиринговой
или расчётной записи прямо отражается на~движении денег.

Главная сложность~--- синхронизация состояния между центрами обработки данных
(ЦОД). Журнал транзакций, таблицы BIN, лимиты для \textit{stand-in} и~ключи HSM
должны оставаться согласованными в~обоих центрах. Для авторизации важнее быстро
восстановить обработку; для клиринга и~расчёта~--- не потерять финансово
значимую запись.

\practicenote{%
  Шардинг по~диапазонам BIN остаётся простой и~распространённой стратегией
  горизонтального масштабирования
  коммутатора авторизации. BIN однозначно определяет эмитента
  и~карточную сеть, а~значит, маршрут и~набор ключей;
  транзакции с~разными BIN не зависят друг от~друга
  и~могут обрабатываться разными экземплярами коммутатора
  без координации. При необходимости шард переносится
  на~другой сервер или в~другой ЦОД без влияния на~остальные
  диапазоны BIN.
}
\subsection*{Аварийное восстановление}

Георезервирование (размещение ЦОД в~разных географических точках) защищает
от~региональных катастроф: пожара, наводнения, отключения электричества
в~районе. Чем дальше площадки друг от~друга, тем лучше защита от~локальных
событий, но~тем выше задержка синхронной репликации. Расстояние выбирают
по~сочетанию регионального риска, качества каналов связи и~цены
синхронной репликации.

RTO (\textit{Recovery Time Objective})~--- время восстановления после полного
отказа ЦОД. Переключение трафика обязано происходить достаточно быстро, чтобы
локальный сбой не превратился в~массовый поток отказов. Парный показатель~---
RPO (\textit{Recovery Point Objective})~--- задаёт допустимую потерю данных:
сколько последних транзакций может не оказаться на~резервной площадке после
переключения. Для авторизации, чей след досверяется по~журналам сети, допустим
больший RPO; для клиринговой и~расчётной записи он стремится к~нулю. Значения
RTO и~RPO определяются архитектурой и~обязательствами оператора.

\section{Подключение к~платёжным системам}\label{sec:pc-connectivity}

\subsection*{VisaNet}

Подключение к~VisaNet, глобальной сети Visa, определяется правилами Visa
и~региональными эксплуатационными регламентами. В~публичных документах Visa
отдельно описаны использование VisaNet, требования к~сетевым сертификатам
и~правила обработки транзакций. VisaNet разделяет авторизацию в~BASE~I
и~пакетный клиринг в~BASE~II.~\cite{visa-core-rules}

Подключение включает канал связи и~сертификацию под~конкретный интерфейс сети.
Публичные правила Visa указывают, что часть деталей, относящихся
к~безопасности сети и~проприетарной внутренней обработке, в~публичной редакции
не раскрывается.~\cite{visa-core-rules}

\subsection*{Banknet}

Mastercard использует собственную сеть Banknet. Публичные материалы Mastercard
описывают её как самостоятельную систему авторизации, клиринга и~расчёта;
Transaction Processing Rules задают сообщения Authorization Request/0100,
Financial Transaction Request/0200, First Presentment/1240 и~форматы
IPM.~\cite{mc-switching-services,mc-tpr}

Mastercard поддерживает single- и~dual-message сценарии; клиринговая часть
опирается на~IPM и~GCMS (\textit{Global Clearing Management System}~---
клиринговая система Mastercard).~\cite{mc-switching-services,mc-tpr}

\subsection*{ОПКЦ НСПК}

Все внутрироссийские карточные транзакции проходят через НСПК~--- независимо
от~того, выпущена ли карта в~Мире или относится к~международной схеме,
работающей внутри страны через локальную инфраструктуру. Банк России
указывает, что внутрироссийская обработка операций международных платёжных
систем ведётся через НСПК; сама НСПК публично описывает себя как оператора ПС
Мир и~ОПКЦ НСПК.~\cite{cbr-ps,nspk-about-company}

Чтобы работать во~внутренней карточной инфраструктуре, банк подключается
к~ОПКЦ НСПК и~соблюдает правила ПС Мир и~связанные технические документы
НСПК.~\cite{nspk-mir-rules}

\subsection*{СБП}

СБП отличается от~карточного процессинга. Она работает как сервис платёжной
системы Банка России. Банк России публикует
отдельный порядок подключения кредитных организаций. НСПК публикует материалы
по~сценариям оплаты по~QR, кнопке и~NFC (\textit{Near Field Communication}~---
бесконтактная радиосвязь ближнего радиуса), а~также
по~возвратам.~\cite{cbr-sbp-actions,sbp-main,sbp-faq-business}

Процессинг, который поддерживает и~карты, и~СБП, совмещает в~одной
инфраструктуре два разных протокола обмена. Карточная операция передаётся
через ISO~8583: в~dual-message сценарии авторизация выполняется в~реальном
времени, клиринг собирается отдельным пакетом, а~сеть рассчитывает
нетто-позиции; в~single-message сценарии финансовое сообщение совмещает
авторизацию и~клиринг. СБП~--- одношаговый перевод средств между счетами
в~Банке России: одно сообщение в~формате ISO~20022, валовый расчёт каждой
операции на~счетах в~Банке России, без отдельного клирингового цикла.
Совмещение требует разных
форматов сообщений, таймаутов, схем сверки и~расчётных моделей внутри одной
платформы обработки.~\cite{cbr-sbp-actions,sbp-main}

\subsection*{Перегрузка и~управляемая деградация}

Когда карточная сеть не~отвечает, процессинговый центр ограничивает ожидание:
тысячи новых авторизаций продолжают поступать, и~очередь быстро переполнится.
Защиту строят на~таймаутах внешних вызовов и~автоматическом размыкании канала
при массовых сбоях; пороги зависят от~реализации процессинга и~обязательств
перед сетями.

\section{Лицензирование и~сертификация}\label{sec:pc-licensing}

\subsection*{Требования Банка России}

Федеральный закон от~27.06.2011~№~161-ФЗ \enquote{О~национальной платёжной
системе} разделяет роли операционного центра (ОЦ), платёжного клирингового
центра (ПКЦ), расчётного центра (РЦ) и~других операторов услуг платёжной
инфраструктуры (ОУПИ); через эти роли описывается юридический статус
процессингового центра. ОУПИ включаются в~реестр операторов услуг платёжной
инфраструктуры, который ведёт Банк
России.~\cite{fz-161,cbr-ps-glossary}

Для самостоятельного ОУПИ базовый набор требований составляют Положение Банка
России от~17.08.2023~№~821-П, Положение Банка России от~13.01.2025~№~850-П,
ГОСТ~Р~57580.1-2017 с~методикой оценки по~ГОСТ~Р~57580.2-2018 и~Федеральный
закон от~27.07.2006~№~152-ФЗ в~части локализации персональных
данных~\cite{cbr-821p,cbr-850p,gost-57580-1,gost-57580-2,fz-152}.
Для~процессингового центра в~составе кредитной организации к~ним добавляется
Положение Банка России от~30.01.2025~№~851-П: оно задаёт требования
к~противодействию переводам без согласия клиента на~стороне
банка-эмитента~\cite{cbr-851p}. Разбор приведён
в~разделе~\ref{sec:regulation-cbr-provisions}. Для платёжной системы Банка
России, включая ОПКЦ СБП и~ССНП (Сервис срочных переводов платёжной системы
Банка России), действует Положение Банка России от~25.07.2022~№~802-П
\enquote{О~защите информации в~платёжной системе Банка
России}~\cite{cbr-802p}.

Для участия в~национальной платёжной системе \enquote{Мир} процессинговый центр
проходит сертификацию для ПС~Мир; её проводит НСПК как оператор
ПС~Мир~\cite{nspk-mir-rules}. Для участия в~СБП требуется интеграция с~ОПКЦ СБП
на~стороне НСПК~\cite{cbr-sbp-rules}. Эти процедуры дополняют, а~не заменяют
требования Банка России.

\subsection*{Сертификация у~карточных сетей}

Каждая карточная сеть проводит собственную программу сертификации для
процессинговых центров, подключающихся к~её инфраструктуре. Публичные документы
Visa и~Mastercard задают условия участия в~сети: правила использования VisaNet,
требования к~обработке транзакций, составу данных, клирингу
и~расчёту.~\cite{visa-core-rules,mc-rules,mc-tpr}

Сертификация охватывает корректность форматов и~обязательных элементов данных,
криптографические процедуры и~устойчивость к~операционным сбоям. Ни Visa, ни
Mastercard в~открытых редакциях не раскрывают полную матрицу тестов, связанных
с~безопасностью. Приёмочные сценарии хранятся в~закрытых руководствах
и~открываются участнику на~этапе
подключения.~\cite{visa-core-rules,mc-rules}

\subsection*{PCI DSS}

Процессинговый центр обрабатывает PAN и~PIN block; issuer processor
дополнительно проверяет ARQC и~генерирует ARPC. Процессинг целиком попадает
в~CDE (cardholder data environment). PCI SSC (\textit{Payment Card Industry
Security Standards Council}) разделяет cardholder data и~SAD (\textit{sensitive
authentication data}). Для PIN, PIN block и~связанных процедур управления
ключами действует самостоятельный стандарт PCI PIN
Security.~\cite{pci-scoping-guidance,pci-faq-1335,pci-pin-security}

Требования к~подтверждению соответствия зависят от~роли и~категории поставщика
услуг в~программе конкретного бренда; одной принадлежности к~Level~1
недостаточно для выбора процедуры. Visa для Level~1 service providers
и~VisaNet processors требует ROC (\textit{Report on Compliance}~--- отчёт
о~соответствии),
подготовленный QSA (\textit{Qualified Security Assessor}~--- аккредитованный
аудитор PCI DSS). Mastercard для регистрируемых поставщиков услуг требует
ежегодного подтверждения соответствия PCI DSS, а~для применимых категорий~---
ежеквартальных сканов ASV (Approved Scanning Vendor~--- одобренный поставщик
внешнего сканирования уязвимостей).~\cite{visa-pci-validation-best-practice,mc-service-provider-categories-pci,pci-faq-1234-asv}

В~CDE процессингового центра попадают коммутатор авторизации, пул HSM,
клиринговый модуль, журнал транзакций и~резервные копии~--- через них проходят
PAN и~SAD. Сетевая сегментация, токенизация PAN на~входе шлюза
и~суррогатные идентификаторы в~аналитических витринах сокращают число систем,
попадающих в~аудит, но~не~выводят процессинг из~CDE целиком (разбор сегментации
и~scoping в~разделе~\ref{sec:pci-cde}).

PCI PIN Security требует обмена ключами PIN-обработки в~формате key block
(ANSI~X9.143, исторически TR-31), запрещает обработку PIN в~открытом виде
вне~HSM и~задаёт key ceremony
(см.~раздел~\ref{sec:crypto-key-ceremony})~\cite{pci-pin-security}.
Процессинговый центр в~терминологии PCI~DSS~--- service provider; Mastercard
относит всех Third Party Processor (TPP) к~Level~1 безусловно, Visa~--- любого,
кто хранит, обрабатывает или передаёт более 300~000 транзакций в~год, а~также
всех VisaNet processor. Процессинг с~таким объёмом или категорией валидируется
как Level~1 через ROC. SAQ~D for Service Providers доступен только
провайдерам ниже порога объёма, не~подпадающим под~автоматическое отнесение
к~Level~1
по~категории~\cite{mc-service-provider-categories-pci,visa-pci-validation-best-practice,pci-saq}.

Для российской финансовой организации PCI DSS действует параллельно
с~ГОСТ~Р~57580.1: оба свода обязательны одновременно и~оцениваются независимыми
циклами (см.~раздел~\ref{sec:pci-russian}).
